%% kapa: exam=1 round=26 year=2015 period=2 form=A subject=accounting subject_ko=회계학 seq=18 number=58
%% source: Q-Net 감정평가사 기출문제 (cst00309) artlSeq=5138001 file=extracted/1008/5138001/제26회_감정평가사 1차_기출/제26회_감정평가사 1차_기출/2교시_감평_A형.pdf
%% conversion: kordoc 4.24.51 -> tools/gen1.py (auto); review=proofread-sonnet; flags=none
\begin{question}{58}
‘농림어업’에 관한 회계처리로 옳지 않은 것은?
\begin{choices}
  \item 생물자산은 최초 인식시점과 매 보고기간말에 공정가치에서 추정 매각부대원가를 차감한 금액(순공정가치)으로 측정하여야 한다. 다만, 공정가치를 신뢰성 있게 측정할 수 없는 경우는 제외한다.
  \item 생물자산에서 수확된 수확물은 수확시점에 순공정가치로 측정하여야 한다.
  \item 생물자산을 최초 인식시점에 순공정가치로 인식하여 발생하는 평가손익과 생물자산의 순공정가치 변동으로 발생하는 평가손익은 발생한 기간의 당기손익에 반영한다.
  \item 수확물을 최초 인식시점에 순공정가치로 인식하여 발생하는 평가손익은 발생한 기간의 기타포괄손익에 반영한다.
  \item 순공정가치로 측정하는 생물자산과 관련된 정부보조금에 다른 조건이 없는 경우에는 이를 수취할 수 있게 되는 시점에만 당기손익으로 인식한다.
\end{choices}
\end{question}
\answer{4}
